\documentclass[11pt]{article}
\usepackage[a4paper,margin=26mm]{geometry}
\usepackage{amsmath,amssymb}
\usepackage[hidelinks]{hyperref}
\setlength{\emergencystretch}{3em}
\title{Conjecture 00000001216: the devil wins on the three-row cylinder}
\author{lizaixi01, with OpenAI Codex assistance}
\date{23 October 2026}
\begin{document}
\maketitle
\section{Board and rules}
The source starts with an angel on $\mathbf Z^2$ and restricts the board
to $\mathbf Z\times[0,2]$. Its lattice points are exactly the three-row
board $S=\mathbf Z\times\{0,1,2\}$. Start at $(0,0)$. We disprove its
assertion that the power-two angel has a winning strategy on this board.
The separate half-plane assertion is not needed for this refutation.

Use Conway's original rules, \href{https://library.slmath.org/books/Book29/files/conway.pdf}
{\emph{The Angel Problem}}, \emph{Games of No Chance}, 1996, p.3:
the devil may eat ``any square of the board whatsoever,'' and the angel
can jump to an uneaten square within its power in the king metric.
In particular the occupied square may be eaten; this forbids future
landings there but does not end the game immediately. Our proof gives
the angel the first move each round, before the devil eats one square.

We strengthen the angel further: it may change to any of the three
rows, and may even stay put, provided its horizontal displacement is
at most two and its landing square was not eaten on an earlier turn.
This includes all usual king-distance-two moves. It also includes all
vertical moves if the three rows are identified cyclically, so the
same obstruction covers that interpretation of ``cylinder.''

\section{An explicit schedule on the infinite board}
The devil uses a fixed schedule, independent of every angel move.
First eat all three squares of each of these four columns, in order:
\[
 -31,\quad -30,\quad 30,\quad 31.
\]
This takes twelve devil turns. In the first twelve angel moves its
horizontal coordinate has absolute value at most $24$, so it has not
passed either double wall when the walls finish. Next eat the interior
columns $-29,-28,\ldots,29$, each in row order $0,1,2$.
There are $59\cdot3=177$ interior squares. Thus the walls and interior
are all eaten after $12+177=189$ devil turns.

For complete precision, define the infinite injective column sequence
\[
 c(k)=\begin{cases}
 -31,-30,30,31&k=0,1,2,3\text{ respectively},\\
 k-33&4\le k<63,\\
 k-31&k\ge63.
 \end{cases}
\]
On zero-based devil turn $n$, eat
$d(n)=(c(\lfloor n/3\rfloor),n\bmod3)$.
All these squares are distinct, including the unused exterior tail.
Consequently this is a legal schedule eating exactly one new square
per turn, not an assumption that the board was initially finite.

\newpage
\section{Why every response loses in bounded time}
Write $a(0)=(0,0)$ and let $a(n+1)$ be the position chosen by the angel
on round $n$, before the devil eats $d(n)$. Legal landing requires
\[
 |a(n+1)_x-a(n)_x|\le2,\qquad
 a(n+1)\ne d(m)\quad(0\le m<n).
\]
The strict inequality $m<n$ is essential: the square eaten immediately
after the current angel move is not incorrectly prohibited in advance.

Induction gives $|a(n)_x|\le2n$ during the initial phase.
In particular $a(12)_x\in[-24,24]\subseteq[-29,29]$.
After the walls finish, suppose the angel lies in $[-29,29]$.
Its next horizontal coordinate lies in $[-31,31]$, and all squares
in the four wall columns $-31,-30,30,31$ are already eaten.
The next coordinate must therefore again lie in $[-29,29]$.
This proves confinement from move twelve onward, using the actual
horizontal move bound and the earlier burns on the infinite board.

Every interior square $(x,y)$ occurs among $d(0),\ldots,d(188)$:
its column index is $k=x+33\in\{4,\ldots,62\}$ and its turn is
$m=3k+y<189$. If the angel could make its 190th move, confinement
would place $a(190)$ in the interior, while legality would prohibit
all 189 earlier burned squares. This is a contradiction.

Thus the devil prevents the 190th angel move after at most 189 burns.
The proof applies to every sequence of legal responses; it imposes
no memory, locality, determinism or computability constraint on them.
In particular no full-history winning angel strategy exists, and the
claimed restricted $7\times7$ strategy cannot exist either.

\section{Formal objects and scope}
Lean uses the actual infinite type $\mathbf Z\times\mathrm{Fin}(3)$.
\texttt{burn\_injective} proves the schedule never repeats a square.
\texttt{LegalPrefix} records exactly the first prescribed number of
angel moves, with the timing above. The bounded induction
\texttt{trapped\_after\_walls} derives the finite interval; it is not
an input assumption. \texttt{no\_legal\_190\_moves} proves finite
failure and \texttt{no\_infinite\_legal\_play} follows.

\texttt{AngelStrategy} permits an arbitrary function of the complete
finite angel and devil histories. \texttt{Follows} records its response
rule. The explicit schedule defeats every such strategy. The final
theorem negates even the necessary condition of surviving all legal
oblivious devil schedules, which any strategy against an adaptive
devil would have to satisfy. No game determinacy theorem is assumed.

The given start $(0,0)$ suffices for the claimed standard game; we do
not separately assert a formal translation theorem for other starts.
Lean 4.33.0 and Mathlib commit
\texttt{db584cd6d46c92f209a44c0f1c829460d327499d} are pinned.
A new root build, fourteen transitive axiom audits and full official
\texttt{leanchecker --fresh Main} replay accompany the source.
No \texttt{sorry}, \texttt{native\_decide} or additional axiom is used.
OpenAI Codex assisted, with independent source/report review by a
separate supervising Codex agent; this is AI review, not human
certification. Organizer acceptance is pending.
\end{document}
